\section*{Capítulo \ref{ch:b3:seismology-imaging} --- Sismología e imagen de la Tierra}

\begin{solution}{exo:b3:seismology-imaging:1}
(a) $\sin i_0 = pv_0/R = 500\times5.80/6371 = 0.4552$, así que $i_0 = \ang{27.1}$.
(b) A \qty{771}{km}, $r = \qty{5600}{km}$ y $r/v = 5600/11.066 =
\qty{506.1}{s.rad^{-1}}$; a \qty{871}{km}, $r = \qty{5500}{km}$ y $r/v =
5500/11.245 = \qty{489.1}{s.rad^{-1}}$. El rayo alcanza su punto de retorno donde $r/v = p = 500$:
$z_t = 771 + 100\times(506.1 - 500)/(506.1 - 489.1) = 771 + 36 =
\qty{807}{km}$.
\end{solution}

\begin{solution}{exo:b3:seismology-imaging:2}
(a) $(455.5 - 369.6)/10 = \qty{8.59}{s}$ por grado, y
$8.59\times180/\pi = \qty{492.2}{s.rad^{-1}}$. (b) Un grado de arco es
$6371\pi/180 = \qty{111.19}{km}$, así que el frente de onda barre la superficie a
$111.19/8.59 = \qty{12.94}{km.s^{-1}}$, mucho más rápido que cualquier roca en la
superficie: el rayo llega casi vertical. (c) En el punto de retorno $v = r_t/p =
5450/492.2 = \qty{11.07}{km.s^{-1}}$, una velocidad de onda P del manto inferior.
\end{solution}

\begin{solution}{exo:b3:seismology-imaging:3}
(a) $\mu_1 = \rho_1\beta_1^2 = 2700\times3500^2 = \qty{3.31e10}{Pa}$
(\qty{33.1}{GPa}); $\mu_2 = 3300\times4500^2 = \qty{6.68e10}{Pa}$
(\qty{66.8}{GPa}). (b) $\lambda = cT = 3.79\times20 = \qty{75.8}{km}$, 2.17
veces el espesor de la capa: la onda penetra bien en el semiespacio, lo que explica
por qué su velocidad ya supera $\beta_1$. (c) $4000/3.384 = \qty{1182}{s}$ y
$4000/4.333 = \qty{923}{s}$: las ondas de 20~s llegan \qty{259}{s}, o
\qty{4.3}{min}, después de las ondas de 100~s.
\end{solution}

\begin{solution}{exo:b3:seismology-imaging:4}
(a) ${}_0S_2$: $1/(53.89\times60) = \qty{3.093e-4}{Hz} = \qty{0.3093}{mHz}$;
${}_0T_2$: $1/(44.01\times60) = \qty{0.3787}{mHz}$. (b) $1440/53.89 = 26.7$
y $1440/44.01 = 32.7$ oscilaciones al día. (c) La diferencia es
$\qty{0.0694}{mHz}$, cuyo inverso es \qty{14400}{s}: al menos \qty{4.0}{h} de
registro, en la práctica mucho más, ya que los picos se ensanchan por la atenuación y
se desdoblan.
\end{solution}

\begin{solution}{exo:b3:seismology-imaging:5}
(a) La cuerda entre dos puntos de la esfera separados por el ángulo $\Delta$
en el centro tiene la longitud $2R\sin(\Delta/2)$, recorrida a $v$.
(b) A \ang{30}: $2\times6371\times\sin\ang{15}/10.0 = \qty{329.8}{s}$, frente a
\qty{369.6}{s}; a \ang{90}: $2\times6371\times\sin\ang{45}/10.0 =
\qty{901.0}{s}$, frente a \qty{779.7}{s}. (c) El rayo hasta \ang{30} retorna en el
manto superior, donde la velocidad real es de unos 8 a \qty{10}{km.s^{-1}}, por debajo
de la de \qty{10}{km.s^{-1}} uniforme; el rayo hasta \ang{90} pasa la mayor parte de su trayecto en el
manto inferior, a 11 a \qty{13.7}{km.s^{-1}}, y se curva hacia dentro, de modo que la
Tierra real es más rápida.
\end{solution}

\begin{solution}{exo:b3:seismology-imaging:6}
(a) Elevando al cuadrado, $(2 - \xi)^4 = 16(1 - \xi/3)(1 - \xi)$. Desarrollando:
$16 - 32\xi + 24\xi^2 - 8\xi^3 + \xi^4 = 16 - \tfrac{64}{3}\xi + \tfrac{16}{3}\xi^2$,
así que $\xi^4 - 8\xi^3 + \tfrac{56}{3}\xi^2 - \tfrac{32}{3}\xi = 0$, es decir,
$\xi\,(3\xi^3 - 24\xi^2 + 56\xi - 32) = 0$, y
$3\xi^3 - 24\xi^2 + 56\xi - 32 = (\xi - 4)(3\xi^2 - 12\xi + 8)$. La raíz
$\xi = 0$ no es una onda. (b) Las raíces de la cuadrática son $2 \pm 2/\sqrt3$,
es decir 3.155 y 0.8453; solo $\xi = 0.8453$ está en $(0, 1)$ (y $\xi = 4$
no), y satisface la ecuación sin elevar al cuadrado, siendo ambos miembros $4/3$.
Así que $c/\beta = \sqrt{0.8453} = 0.9194$. (c) $c = 0.9194\times3.46 =
\qty{3.181}{km.s^{-1}}$; sobre \qty{1000}{km}, $314.4 - 289.0 = \qty{25.3}{s}$
después de la onda S.
\end{solution}

\begin{solution}{exo:b3:seismology-imaging:7}
$\mathsf G$ tiene las filas $(1, x_i)$ y $\mathsf C_d^{-1} = \operatorname{diag}
(100, 100, 25, 25)$ (\unit{s^{-2}}). Las ecuaciones normales son
\[
  \begin{pmatrix} 250 & 23750 \\ 23750 & \num{2812500} \end{pmatrix}
  \begin{pmatrix} a \\ b \end{pmatrix}
  = \begin{pmatrix} 3665 \\ 418125 \end{pmatrix},
\]
con determinante \num{1.390625e8}. Por tanto, $b = (250\times418125 -
23750\times3665)/\num{1.390625e8} = \qty{0.12575}{s.km^{-1}}$ y
$a = (2812500\times3665 - 23750\times418125)/\num{1.390625e8} =
\qty{2.713}{s}$. La covarianza es $\mathsf N^{-1}$: $\sigma_b =
\sqrt{250/\num{1.390625e8}} = \qty{0.00134}{s.km^{-1}}$, $\sigma_a =
\sqrt{2812500/\num{1.390625e8}} = \qty{0.142}{s}$. La velocidad es
$1/b = \qty{7.95}{km.s^{-1}}$, con el error $\sigma_b/b^2 =
\qty{0.08}{km.s^{-1}}$: una velocidad del manto más superficial. Los residuos son $-0.001$,
$0.011$, $-0.076$ y $\qty{0.036}{s}$, y $\chi^2 = 0.19$ para $4 - 2 = 2$
grados de libertad: menor de lo esperado, así que los errores indicados son, si acaso,
generosos; el ajuste es aceptable.
\end{solution}

\begin{solution}{exo:b3:seismology-imaging:8}
(a) $\mathsf G = (30\;\;40)$ es una única fila: $\mathsf G^{\mathsf T}\mathsf G$ es
de $2\times2$ y de rango uno, por tanto singular; cualquier modelo con $30m_1 + 40m_2 = 0$ puede
añadirse sin cambiar el dato. (b) $\mathsf G\mathsf G^{\mathsf T} =
900 + 1600 = 2500$, así que $\mathbf m = (30, 40)\times0.50/2500 =
(0.006, 0.008)\,\unit{s.km^{-1}}$. (c) $\mathsf R = \frac{1}{2500}
\begin{pmatrix} 900 & 1200 \\ 1200 & 1600 \end{pmatrix} =
\begin{pmatrix} 0.36 & 0.48 \\ 0.48 & 0.64 \end{pmatrix}$, de traza 1: un dato
resuelve una combinación de parámetros. (d) $\mathsf R\,(0.01, 0) =
(0.0036, 0.0048)\,\unit{s.km^{-1}}$: la anomalía se reparte a lo largo del rayo, más
en la celda donde el rayo es más largo.
\end{solution}

\begin{solution}{exo:b3:seismology-imaging:9}
(a) Linealizando: $\delta t = -(\delta v/v)L/v = -0.03\times200/9.0 =
\qty{-0.667}{s}$; exactamente, $L/(1.03v) - L/v = 200/9.27 - 200/9.0 =
\qty{-0.647}{s}$. (b) El trayecto en la losa es $100/\cos\ang{40} =
\qty{130.5}{km}$: $-0.03\times130.5/9.0 = \qty{-0.435}{s}$ linealizado y
$\qty{-0.422}{s}$ exacto. (c) El valor exacto es el linealizado dividido por
$1.03$: un error de \qty{3}{\percent}, del orden de $\delta v/v$ en sí, como
cabe esperar de una teoría de primer orden.
\end{solution}

\begin{solution}{exo:b3:seismology-imaging:10}
(a) El sobretono existe mientras su parte en el semiespacio decae, $c < \beta_2$; en su
frecuencia de corte $c = \beta_2$, así que $q_2 = 0$ y $\tan(kHq_1) = 0$, con $kHq_1$ en
$[n\pi, n\pi + \pi/2)$: $kHq_1 = n\pi$. (b) Con $k = \omega/\beta_2$ y
$q_1 = \sqrt{\beta_2^2/\beta_1^2 - 1} = \sqrt{\beta_2^2 - \beta_1^2}/\beta_1$,
$\omega_nH\sqrt{\beta_2^2 - \beta_1^2}/(\beta_1\beta_2) = n\pi$, que es la
fórmula. (c) $\sqrt{4.5^2 - 3.5^2} = \sqrt8 = 2.828$; $\omega_1 =
\pi\times3.5\times4.5/(35\times2.828) = \qty{0.4998}{rad.s^{-1}}$, un periodo de
$2\pi/0.4998 = \qty{12.57}{s}$: a periodos más largos solo existe el modo fundamental.
existe. (d) Para $n = 0$ la frecuencia de corte es cero: el modo fundamental
existe a todos los periodos. A periodos cortos $kH \to \infty$ mientras $kHq_1$ se mantiene
por debajo de $\pi/2$, así que $q_1 \to 0$ y $c \to \beta_1$.
\end{solution}

\begin{solution}{exo:b3:seismology-imaging:11}
(a) $\mathsf G^{\mathsf T}(\mathsf G\mathsf G^{\mathsf T} + \varepsilon^2\mathsf I)
= (\mathsf G^{\mathsf T}\mathsf G + \varepsilon^2\mathsf I)\mathsf G^{\mathsf T}$;
se multiplica por la izquierda por $(\mathsf G^{\mathsf T}\mathsf G + \varepsilon^2
\mathsf I)^{-1}$ y por la derecha por $(\mathsf G\mathsf G^{\mathsf T} +
\varepsilon^2\mathsf I)^{-1}$, ambas invertibles para $\varepsilon > 0$.
(b) Aquí $\mathsf G\mathsf G^{\mathsf T} = \|\mathbf g\|^2 = 2500$, un número,
por tanto $\mathsf R_\varepsilon = \mathsf G^{\mathsf T}\mathsf G/(2500 + \varepsilon^2)$,
de traza $2500/(2500 + \varepsilon^2)$. (c) El dato predicho es
$\mathbf g\cdot\mathbf m_\varepsilon = 0.50\times2500/(2500 + \varepsilon^2)$
y el residuo $0.50\,\varepsilon^2/(2500 + \varepsilon^2)$. Para
$\varepsilon^2 = 0$: traza 1, residuo 0; para 2500: traza 0.5, residuo
\qty{0.25}{s}; para \num{10000}: traza 0.2, residuo \qty{0.40}{s}. Un mayor amortiguamiento
da lugar a un modelo más pequeño, una resolución peor y un desajuste mayor.
\end{solution}

\begin{solution}{exo:b3:seismology-imaging:12}
(a) Por el ejercicio 5, $T = (2R/v)\sin(\Delta/2)$, así que $\mathrm dT/\mathrm d\Delta
= (R/v)\cos(\Delta/2)$. Una cuerda que forma el ángulo $\Delta$ en el centro es la
base de un triángulo isósceles de ángulo en el vértice $\Delta$, por lo que corta al radio
en la fuente con $i_0 = (\pi - \Delta)/2$, y $p = R\sin i_0/v =
(R/v)\cos(\Delta/2)$: lo mismo. (b) El punto de retorno es el punto medio de la
cuerda, a la distancia $R\cos(\Delta_1/2)$ del centro; equivalentemente
$r_1 = p_1v$. (c) Aquí $p(\Delta)/p_1 = \cos(\Delta/2)/\cos(\Delta_1/2)$, y
la fórmula da $r_1 = R\exp\bigl(-\ln(1/\cos(\Delta_1/2))\bigr) =
R\cos(\Delta_1/2)$, como en (b); la velocidad es $r_1/p_1 = R\cos(\Delta_1/2)\big/
\bigl((R/v)\cos(\Delta_1/2)\bigr) = v$.
\end{solution}

\begin{solution}{pb:b3:seismology-imaging:1}
\textbf{1.} $i_0 = \arcsin 0.8 = \ang{53.13}$ y $p = 0.8/5.0 =
\qty{0.16}{s.km^{-1}}$.

\textbf{2.} El rayo alcanza su punto de retorno donde $v = 1/p = \qty{6.25}{km.s^{-1}}$, es decir, en
$z = (6.25 - 5.0)/0.05 = \qty{25}{km}$.

\textbf{3.} Derivando $\sin i = pv$ a lo largo del rayo, $\cos i\,\mathrm di/
\mathrm ds = pg\,\mathrm dz/\mathrm ds = pg\cos i$, así que $\mathrm di/\mathrm ds =
pg$: la dirección gira a ritmo constante, con la curvatura de un círculo de radio
$1/(pg) = 1/(0.16\times0.05) = \qty{125}{km}$. En el punto de retorno, el rayo es
horizontal a \qty{25}{km} de profundidad, y el centro se encuentra verticalmente a \qty{125}{km}
por encima de él: en $z = \qty{-100}{km}$, es decir, a $v_0/g$ por encima de la superficie, donde la
velocidad extrapolada se anularía.

\textbf{4.} En la fuente, el radio del círculo es perpendicular al rayo,
por lo que forma el ángulo $i_0$ con la horizontal, y el centro queda
a $125\cos i_0$ de distancia horizontal; por simetría
$X = 2\times125\times\cos i_0 = 2\times125\times0.6 = \qty{150}{km}$, o
$X = 2\sqrt{1 - p^2v_0^2}/(pg)$.

\textbf{5.} $\mathrm dT = \mathrm ds/v$ con $\mathrm ds = \mathrm di/(pg)$ y
$v = \sin i/p$: $\mathrm dT = \mathrm di/(g\sin i)$. Entonces
\[
  T = \frac2g\int_{i_0}^{\pi/2}\frac{\mathrm di}{\sin i} = \frac2g\ln\cot\frac{i_0}2 =
  \frac2g\operatorname{arccosh}\frac1{pv_0} .
\]
Aquí $\cot(\ang{26.57}) = 2$, así que
$T = 40\ln2 = \qty{27.73}{s}$.

\textbf{6.} Para $p = 0.15$: $pv_0 = 0.75$, $\sqrt{1 - 0.5625} = 0.6614$,
$X = 2\times0.6614/0.0075 = \qty{176.4}{km}$ y $T = 40\operatorname{arccosh}
(1.3333) = 40\times0.7954 = \qty{31.81}{s}$. Entonces $\Delta T/\Delta X =
4.08/26.4 = \qty{0.155}{s.km^{-1}}$, entre 0.15 y 0.16, como
$\mathrm dT/\mathrm dX = p$ lo exige.

\textbf{7.} $\mu_1 = 2700\times3500^2 = \qty{33.1}{GPa}$ y
$\mu_2 = 3300\times4500^2 = \qty{66.8}{GPa}$.

\textbf{8.} $k = 2\pi/(cT) = 2\pi/(3.99\times25) = \qty{0.06299}{rad.km^{-1}}$;
$q_1 = \sqrt{(3.99/3.5)^2 - 1} = 0.5474$; $q_2 = \sqrt{1 - (3.99/4.5)^2} = 0.4624$.

\textbf{9.} $\tan(kHq_1) = 66.825\times0.4624/(33.075\times0.5474) = 1.707$; el
modo fundamental toma $kHq_1 = \arctan1.707 = 1.041$, así que
$H = 1.041/(0.06299\times0.5474) = \qty{30.2}{km}$.

\textbf{10.} A \qty{40}{s}: $k = 2\pi/170 = 0.03696$, $q_1 = 0.6888$,
$q_2 = 0.3287$, $\tan(kHq_1) = 0.964$, $kHq_1 = 0.7671$ y
$H = 0.7671/(0.03696\times0.6888) = \qty{30.1}{km}$. Los dos resultados coinciden dentro de la
precisión de las velocidades de fase: la corteza tiene un espesor de unos \qty{30}{km}.

\textbf{11.} $\omega = 0.2513$ y $\qty{0.1571}{rad.s^{-1}}$, $k = 0.06299$ y
$\qty{0.03696}{rad.km^{-1}}$: $U \approx 0.0942/0.02603 = \qty{3.62}{km.s^{-1}}$,
menor que ambas velocidades de fase, como debe ser cuando $c$ crece con el
periodo.

\textbf{12.} $\lambda = 3.99\times25 = \qty{99.8}{km}$ y $4.25\times40 =
\qty{170}{km}$. El modo fundamental muestrea hasta una fracción de su
longitud de onda: la onda de 25 s se sitúa sobre todo en la corteza, la de 40 s llega bien
dentro del manto, así que su velocidad está más cerca de $\beta_2$. Dos periodos que fijan el
mismo $H$ es una comprobación de que el modelo de dos capas es adecuado.

\textbf{13.} Con las columnas de las celdas 1 a 4 y las filas A a F, en unidades
de \qty{1000}{km} ($\sqrt2\times100/1000 = 0.1414$):
\[
  \mathsf G = \begin{pmatrix}
  0.1 & 0.1 & 0 & 0 \\ 0 & 0 & 0.1 & 0.1 \\ 0.1 & 0 & 0.1 & 0 \\
  0 & 0.1 & 0 & 0.1 \\ 0.1414 & 0 & 0 & 0.1414 \\ 0 & 0.1414 & 0.1414 & 0
  \end{pmatrix}.
\]

\textbf{14.} $\mathsf C_d^{-1} = \operatorname{diag}(400, 400, 400, 400, 100,
100)$. Un rayo de fila o columna aporta $0.1^2\times400 = 4$, un rayo diagonal
$0.1414^2\times100 = 2$. La celda 1 es atravesada por A y C (4 cada uno) y por E (2): 10. Las celdas
1 y 2 comparten el rayo A: 4; las celdas 1 y 3 comparten C: 4; las celdas 1 y 4 comparten E: 2; y
así sucesivamente por simetría, lo que da las cuatro filas indicadas.

\textbf{15.} Un rayo de fila o columna aporta $0.1\times400\,\delta t = 40\,
\delta t$ a cada una de sus celdas, y un rayo diagonal $0.1414\times100\,\delta t =
14.14\,\delta t$. Celda 1: $40(0.16 - 0.36) + 14.14\times(-0.15) = -10.12$; celda
2: $40(0.16 + 0.11) + 14.14\times(-0.27) = 6.98$; celda 3: $40(-0.48 - 0.36) +
14.14\times(-0.27) = -37.42$; celda 4: $40(-0.48 + 0.11) + 14.14\times(-0.15) =
-16.92$.

\textbf{16.} Para $(1,1,1,1)$, cada fila de $\mathsf N$ suma 20. Para
$(1,-1,-1,1)$, la fila 1 da $10 - 4 - 4 + 2 = 4$ y las demás siguen el mismo
patrón. Para $(1,1,-1,-1)$, $10 + 4 - 4 - 2 = 8$, y para $(1,-1,1,-1)$,
$10 - 4 + 4 - 2 = 8$. Los vectores son ortogonales con norma 2, así que
$\mathsf N^{-1} = \sum_k \mathbf u_k\mathbf u_k^{\mathsf T}/(4\lambda_k)$; su
elemento $(1,1)$ es $\tfrac14(\tfrac1{20} + \tfrac14 + \tfrac18 + \tfrac18) =
\tfrac{11}{80}$, su elemento $(1,2)$ es $\tfrac14(\tfrac1{20} - \tfrac14 + \tfrac18
- \tfrac18) = -\tfrac4{80}$, su elemento $(1,4)$ es $\tfrac14(\tfrac1{20} +
\tfrac14 - \tfrac18 - \tfrac18) = \tfrac1{80}$, y así sucesivamente.

\textbf{17.} $\hat{\mathbf m} = \mathsf N^{-1}\mathbf b$:
$\hat m_1 = (11\times(-10.121) - 4\times6.982 - 4\times(-37.418) - 16.921)/80
= -0.081$; del mismo modo $\hat m_2 = 1.844$, $\hat m_3 = -3.706$ y
$\hat m_4 = \qty{-0.931}{ms.km^{-1}}$.

\textbf{18.} $\delta v/v_0 = -v_0\,\delta s = -8.0\times10^{-3}\,\hat m$, en tanto por
ciento $-0.8\,\hat m$: celda 1 $+0.07$, celda 2 $-1.48$, celda 3 $+2.96$, celda
4 $\qty{+0.75}{\percent}$.

\textbf{19.} Todo elemento diagonal de $\mathsf N^{-1}$ vale $11/80$: cada celda
tiene el error estándar $\sqrt{0.1375} = \qty{0.371}{ms.km^{-1}}$, es decir,
$\qty{0.30}{\percent}$. La correlación de las celdas 1 y 2 es $-4/11 = -0.36$:
comparten el rayo A, y una lentitud asignada a una se puede compensar con la otra.
Las celdas 1 y 4, que solo comparten la diagonal E, de menor precisión, están
débilmente y positivamente correlacionadas, $+1/11 = +0.09$.

\textbf{20.} Los residuos predichos son $0.176$, $-0.464$, $-0.379$, $0.091$,
$-0.143$ y $\qty{-0.263}{s}$; los residuos tras el ajuste son $-0.016$,
$-0.016$, $0.019$, $0.019$, $-0.007$ y $\qty{-0.007}{s}$, y
$\chi^2 = 0.50$ para $6 - 4 = 2$ grados de libertad: coherente con los errores
indicados, si acaso algo pequeños.

\textbf{21.} Los autovectores son los de $\mathsf N$ y los autovalores son
$\lambda/(\lambda + 4)$: $20/24 = 5/6$, $4/8 = 1/2$, $8/12 = 2/3$ dos veces. Como
en la pregunta 16, $R_{11} = \tfrac14(\tfrac56 + \tfrac12 + \tfrac23 + \tfrac23) =
\tfrac23$; $R_{12} = \tfrac14(\tfrac56 - \tfrac12 + \tfrac23 - \tfrac23) =
\tfrac1{12}$; $R_{13} = \tfrac14(\tfrac56 - \tfrac12 - \tfrac23 + \tfrac23) =
\tfrac1{12}$; $R_{14} = \tfrac14(\tfrac56 + \tfrac12 - \tfrac23 - \tfrac23) = 0$.

\textbf{22.} Proyéctese $\mathbf b$ sobre los autovectores: $\mathbf b\cdot\mathbf u_k/4
= -14.37$, $0.849$, $12.80$ y $-9.40$ para los autovalores 20, 4, 8 y 8.
La estimación amortiguada los divide por $\lambda + 4$ en lugar de por $\lambda$; para
la celda 3 (componentes $1, -1, -1, 1$),
$m_3 = -14.37/24 - 0.849/8 - 12.80/12 - 9.40/12 = \qty{-2.555}{ms.km^{-1}}$,
una anomalía de velocidad de $\qty{+2.04}{\percent}$. El amortiguamiento lleva cada componente
hacia cero, sobre todo las peor restringidas.

\textbf{23.} El modelo sin amortiguar, que ajusta los datos dentro de sus errores
($\chi^2 = 0.50$ para 2 grados de libertad), mientras que el amortiguado los
desajusta notablemente ($15.8$). El amortiguamiento solo es necesario cuando
$\mathsf N$ es singular o casi lo es; aquí su autovalor menor, 4, no es pequeño, la
estimación es estable y sus errores son moderados.

\textbf{24.} La fila 3 es $(1/12, 0, 2/3, 1/12)$: la estimación amortiguada de la
celda 3 es dos tercios de su valor verdadero más una doceava parte de las celdas 1 y 4,
sus vecinas a lo largo del rayo C y del rayo B, y nada de la celda 2.

\textbf{25.} La celda 3, la celda inferior de la izquierda, es rápida en unos
\qty{3}{\percent}; la celda 2 es lenta. Una losa fría y densa de litosfera oceánica
que se hunde bajo el continente es más rápida que el manto circundante, así que
la celda rápida, profunda y en el lado oceánico de K, es la losa; la celda lenta
situada encima, al otro lado, se encuentra donde la cuña del manto está caliente.

\textbf{26.} A partir de las preguntas 18 y 19, la anomalía de velocidad P de la celda de la losa es
\textbf{\qty[text-series-to-math=true, uncertainty-mode=separate]{3.0 \pm 0.3}{\percent}}.
\end{solution}
